\documentclass[runningheads]{llncs}

\usepackage{amsmath,amssymb}
\usepackage{graphicx}
\usepackage{booktabs}
\usepackage{tabularx}
\usepackage{array}
\usepackage{url}
\usepackage{hyperref}
\usepackage{enumitem}

\begin{document}

\title{Applying MESA Benchmark Reviews to Requirement Decisions: An Executable Method}
\titlerunning{Applying MESA Reviews to Requirement Decisions}

\author{Anonymous Author(s)}
\authorrunning{Anonymous Author(s)}

\institute{Affiliation Withheld for Double-Blind Review\\
\email{contact@example.org}}

\maketitle

\begin{abstract}
\textbf{Context and motivation:} Practitioners who adopt coding agents often treat public benchmark results as evidence that an agent meets a requirement. MESA, a rubric for reviewing AI benchmarks presented in a separate submission, has rules that tie a requirement to the review evidence it needs before benchmark scores are compared with a threshold. Automated audits already find flawed tasks in benchmarks, but whether a review changes a requirement decision beyond an audit has not been tested.

\textbf{Question/problem:} RQ1: How does reviewing a benchmark with MESA or ABC, beyond an automated audit, change the decision on whether a coding agent meets a requirement? RQ2: Does a decision from a reviewed benchmark agree with a direct test of the same requirement?

\textbf{Principal ideas/results:} We apply MESA's rules to a requirement on a coding agent, with records and a decision engine that report input integrity, evidence status, and the decision separately and return every blocking reason. A rating counts only when its check passed, its scope covers the coding agent and tasks, and the evidence the requirement asks for is present; a task with shown broken tests is left out before any score is seen. This paper adds a signed transfer justification, execution provenance for every observation, and a strict and a lenient policy. The Wilson bound shows that at a threshold of 0.70 Accept needs at least nine observations (9 of 9) and, with 20, 19 or 20 resolved, and calculations in Section~\ref{sec:completed} show that with 30 observations and a true resolve rate of 0.75, 89.4\% of decisions remain inconclusive. We prepared an evaluation on SymPy with two coding agents, 20 source tasks from SWE-bench Verified, and 20 target tasks from bug reports fixed after both models' stated training cutoffs, which form the direct test; at most 90 sessions are planned. Claude Sonnet 5.5 makes exploit attempts against the scoring of the source tasks (a patch passed the hidden tests and failed an independent check on 17 of the 20 tasks) and a grading check (patches that change only test running files were graded as resolved in 9 of 9 cases), writes the rating forms in two separate runs, and answers the items of the Agentic Benchmark Checklist (ABC). As the forms stand, MESA withholds its decision under both policies, while the resolve rate alone, an audit only decision, and ABC each give one. No study session has been run.

\textbf{Contribution:} (1)~Executable records and checks for identity, artifact hashes, scope, and execution provenance, and a decision engine that keeps evidence defects, missing information, and sample uncertainty apart; (2)~a case builder that compares the resolve rate alone, audit only decisions, MESA under two policies, and ABC, and sets each against the direct test; (3)~a session validity rule and its count by kind, as a control; and (4)~an evaluation design with 20 source and 20 target tasks per coding agent, exploit attempts, a grading check, two separate rating runs, and a replay that tests the implementation. Session, rating, and replay observations remain future work.

\keywords{Requirements Engineering for AI \and AI Benchmarks \and Acceptance Testing \and Evidence Reuse \and Traceability \and Coding Agents.}
\end{abstract}

\section{Introduction}
\label{sec:intro}

Projects that select a coding agent need evidence that it meets requirements. We use one throughout: ``Working alone on a bug report from the SymPy repository, the coding agent changes the code so that the tests the maintainers later wrote for that bug pass, for at least 70\% of bug reports, within 15 minutes each.'' It is an illustrative requirement written for this study, not one the SymPy maintainers stated. Its form follows the traceability example in the MESA paper, a threshold on a benchmark resolve rate, and the way requirements engineering for AI states quality requirements: a measurable threshold traced to evidence \cite{cleland_huang_traceability_2012,habiba_re4ai_maturity_2024,lewis_quality_model_ml_2026}. It is one acceptance criterion; a project that adopts a coding agent would also weigh review effort and security. The values 70\% and 15 minutes were fixed before any session, and the method works with any values.

We use these terms throughout. A \emph{coding agent} is a model in its tool; the records call it a configuration. An \emph{issue} is a bug report, and a \emph{task} is an issue prepared for an agent with the tests of the real fix hidden. A \emph{session} is one coding agent run on one task; MESA's rating sessions are called rating runs here. A task is \emph{resolved} when every hidden test passes, and the \emph{resolve rate} is the share of tasks resolved. MESA, the rubric for reviewing AI benchmarks, asks 31 evaluation questions, called \emph{criteria}; the \emph{level} is the answer to one criterion, such as Verified or Contradicted. The \emph{decision} is the method's answer about a requirement: Accept, Reject, or Inconclusive, or no decision when the evidence does not allow one. The \emph{direct test} has each coding agent work on bug reports from the repository the requirement is written for, fixed after the models' stated training cutoffs; the same rule gives its decision, and its tasks are the \emph{target tasks}.

A resolve rate supports a requirement only if the tests check the required behavior, the runs were isolated, and the tasks were not exposed in training \cite{kane_validation,jo_wilson_what_does_your_benchmark_measure}. Practitioners still consult leaderboards in place of direct tests of their own requirements \cite{ghosh_evaluation_cards_2026,hou_llm4se_survey_2024,singh_leaderboard_illusion}, while audits of coding benchmarks found tests that accept incomplete patches or reject correct ones, underspecified tasks, and training data exposure \cite{ding_spec_oracle_gaps_2026,openai_swebench_audit_2026,openai_swebench_verified_contamination_2025}. MESA, presented in a separate submission (reference withheld for review), addresses the review side. It maps 280 items from benchmark review frameworks to 30 shared concerns and derives a rubric of eight description facets and 31 evaluation questions, each answered with one of six levels (Section~\ref{sec:foundation}). It also states rules for a requirements traceability matrix: a requirement names the evidence it needs, an evidence status is set from the ratings, and a score decision is made only when the evidence is admissible. This paper applies those rules to a requirement on a coding agent and prepares their evaluation. It does not revalidate MESA and does not use the benchmark ratings reported in the MESA paper as observations.

\emph{Gap 1.} AI agents already audit real benchmarks and find confirmed flaws: BenchGuard \cite{tu_benchguard_2026}, ABA \cite{wang_auto_benchmark_audit_2026}, which on SWE-bench Verified shows that filtering flagged tasks changes rankings, and BenchJack \cite{wang_benchjack_2026}, whose coding agents run exploits. UTBoost \cite{yu_utboost_2025} and SWE-bench+ \cite{aleithan_swebench_plus_2024} show that resolve rates move when flawed tasks are removed, and ABC \cite{zhu_agentic_benchmark_checklist_2025} scores items but gives no rule to accept or reject a result. MESA's own studies (the companion paper, reference withheld) rated benchmarks from files only, so no exploit check, named model, or defective task entered a rating, and its one worked traceability entry uses ratings without observations. None turns a benchmark review into a decision on a stated requirement (threshold, confidence, task set) or says what evidence would be needed to reach one. This gives RQ1.

\emph{Gap 2.} Benchmark results stand in for direct tests, as noted above. SWE-bench Live and SWE-rebench show score gaps on newer tasks for whole benchmarks \cite{zhang_swebench_live_2025,badertdinov_swe_rebench_2025}, and the SWE-Bench Illusion shows signs of memorization \cite{liang_swebench_illusion_2025}. None checks a decision from a reviewed benchmark against a direct test of the same requirement. This gives RQ2.

\subsection{Research Questions}
\begin{itemize}[leftmargin=*]
    \item \textbf{RQ1:} How does reviewing a benchmark with MESA or ABC, beyond an automated audit, change the decision on whether a coding agent meets a requirement?
    \item \textbf{RQ2:} Does a decision from a reviewed benchmark agree with a direct test of the same requirement?
\end{itemize}
RQ1 is answered, for each coding agent, by comparing five decisions on the source tasks: the resolve rate alone, audit only, ABA as a second audit, MESA under two policies, and ABC (Section~\ref{sec:evaluation}), with the tasks each leaves out, the evidence status, the blocking criteria, and the caveats. RQ2 is answered on 20 target tasks per coding agent, which are the direct test, by setting each of these decisions against the direct test's decision. The rating forms, exploit attempts, grading check, and ABC record are produced before any session; they are the inputs to RQ1 and the preliminary validation of the artifacts, and the author's review and signing of them is a step of the method. Session validity is a control of the method, reported but not a research question. The completed calculations (Section~\ref{sec:completed}) give context: they show how many tasks a decision needs.

\subsection{Contributions and Claims}
This paper contributes (1)~versioned case records and rating form records that link a requirement, its evidence conditions, a coding agent, task identifiers, observations, artifact hashes, the MESA criterion with its definition, level, outcome, check kind, scope, detection limits, task defects, and the roles of the author and the models, and a decision engine that keeps input integrity, evidence status, and decision apart and returns every reason; (2)~a case builder that turns sessions, signed rating forms, the ABC record, the requirement, and the transfer justification into cases, one row per method, and sets each decision against the direct test; (3)~a session validity rule applied to every transcript, as a control; and (4)~an evaluation design on SymPy with 20 source and 20 target tasks per coding agent and at most 90 sessions, exploit attempts and a grading check on the source tasks, two separate rating runs, and a replay that tests the implementation. Table~\ref{tab:claims} separates what this paper takes from the MESA paper, what it implements, what has been calculated, and what remains to be observed.

\begin{table}[t]
\caption{Claims and their evidence. Fixture tests check software behavior on constructed inputs; they are not study findings.}
\label{tab:claims}
\centering
\scriptsize
\setlength{\tabcolsep}{3pt}
\begin{tabularx}{\textwidth}{@{}>{\raggedright\arraybackslash}p{4.6cm}>{\raggedright\arraybackslash}p{2.2cm}>{\raggedright\arraybackslash}X@{}}
\toprule
\textbf{Claim} & \textbf{Status} & \textbf{Evidence} \\
\midrule
\multicolumn{3}{@{}l}{\emph{Foundations from the MESA paper}} \\
Shared concerns, the seven criteria B05, B07, B12, B14a, B14b, B28a, and B28b, six levels, and the record fields & Taken as given & MESA rubric, cited by path and hash \\
Evidence status, score decision, removal of defective tasks, and the requirement's task set, configuration, minimum observations, threshold, and confidence & Taken as given & MESA paper; not tested here \\
\midrule
\multicolumn{3}{@{}l}{\emph{Implementation contributions of this paper}} \\
Records link requirement, configuration, tasks, observations, hashes, concern, scope, check, conditions, and roles & Implemented & Code and fixture tests \\
Identity and hash checks precede any capability judgment & Implemented & Fixture tests \\
A \emph{Verified} rating alone does not satisfy a condition & Implemented & Fixture tests \\
Source and target decisions are separate; transfer needs a justification & Implemented & Fixture tests \\
Replay copies keep parent digests & Implemented & Fixture tests \\
A task with shown broken tests is left out of M3 and named; the score is reported with and without it & Implemented & Fixture tests \\
Session ledger enforces phase limits and refuses a 91st session & Implemented & Fixture tests with stub adapters \\
\midrule
\multicolumn{3}{@{}l}{\emph{Completed calculations}} \\
Decision probabilities by sample size & Completed & Exact binomial enumeration (Table~\ref{tab:power_analysis}) \\
Accept needs at least nine observations at 0.70 & Calculated & Wilson bound: 9 of 9 gives 0.701, 8 of 8 gives 0.676 (Section~\ref{sec:planned}); fixture test \\
Effect of assumed review accuracy on false acceptance & Completed under assumptions & Saved simulation outputs (Table~\ref{tab:simulation_far}) \\
\midrule
\multicolumn{3}{@{}l}{\emph{Checks run and observations pending}} \\
SymPy task eligibility; target task identities & Run & Harness runs: 75 of 75 source tasks and the first 20 target candidates eligible; identities confirmed through the GitHub API \\
Agent tool sandboxes (network, reads outside the workspace) & Run & Codex blocks network access from shell commands and writes outside the workspace, not outside reads; its web search and its apps server, which also searched the web, are turned off, after which a check found no web search tool. The Antigravity tool runs every shell command in its sandbox, which blocked network access and reads and writes outside the workspace, as an agent without web tools; a hook refuses its file tools outside the workspace and requests to leave the sandbox, and the run goes on. Codex keeps other tools outside the shell \\
Exploit attempts on the 20 source tasks (B07); grading check (B28a) & Done & Attempts recorded for all 20 tasks; on 17 a patch passed the hidden tests and failed the independent check; patches that change only test running files were graded as resolved in 9 of 9 cases \\
ABC items answered by Claude Sonnet 5.5 and the author's signature & Answered; signature pending & 8 items met, 9 not met, 10 not applicable, 6 unknown; ABC decides on all 20 source tasks, as the resolve rate alone does, with 15 caveats \\
Two rating runs by Claude Sonnet 5.5, their agreement, and the author's signatures & Done; signatures pending & Eight forms per run; same level on 7 of 8 forms, same outcome on 6 of 8, same task defects on all 8; they differ on the outcome of B05 and on B28a \\
Source and target decisions for two coding agents, and their agreement & Pending & Result templates and the case builder prepared \\
Replay on collected records (two changed copies per task) & Pending & Definitions and tooling prepared; a test of the implementation, not a research question \\
Runtime model identity, the time limit, isolation & Pending & Observed only in recorded sessions (Section~\ref{sec:threats}) \\
\bottomrule
\end{tabularx}
\end{table}

\section{Background and Related Work}
\label{sec:related_work}

\subsection{What This Paper Takes from MESA}
\label{sec:foundation}
The MESA rubric has eight description facets and 31 evaluation criteria, each a question about a benchmark. A \emph{rating form} records the \emph{level} of one criterion with its evidence and checks. A level is one of \emph{Claimed}, \emph{Documented}, \emph{Verified}, \emph{Contradicted}, \emph{Undocumented}, and \emph{Not applicable}. \emph{Contradicted} is the only negative level and needs a cited location. Missing information is not a defect. The rubric has no total score. Every rating records its level, its outcome (passed, failed, not performed, or not recorded), the kind of check (reported by the sources, run by the reviewer, both, or none), its scope, its detection limits, any task defects (task, direction, location, and quote), and who rated and who signed.

MESA's traceability rules, which we apply, are these. A requirement states its task set, the model configuration, the minimum number of observations, the threshold, and the confidence level. The evidence status is \emph{Admissible}, \emph{Inquiry Required}, or \emph{Inadmissible}. A score decision (\emph{Accept}, \emph{Reject}, or \emph{Inconclusive}) is made only when the status is Admissible; otherwise there is no decision, and every blocking criterion is listed. Defective tasks are removed before any score is seen, and the score is reported with and without them. Following Kane \cite{kane_validation}, MESA treats a benchmark score as support for a claim only through a chain of scoring, generalization, extrapolation, and decision inferences. In MESA's mapping, B07 belongs to the scoring inference, B12 to generalization, B05, B14a, and B14b to extrapolation, and B28a and B28b to the decision inference. The transfer justification used here is a short argument, signed by the author, that source tasks predict the requirement on the target; it concerns the step from source tasks to the target.

This paper adds three rules of its own to MESA's.
\begin{enumerate}[leftmargin=*]
    \item The decision on the requirement is the decision on the source tasks, taken through a signed transfer justification; without one there is no decision on the requirement.
    \item Every observation carries execution provenance: a session, a log, and the log's hash.
    \item A strict policy requires \emph{Verified} with a passed check for all seven criteria, and a lenient policy requires less (Section~\ref{sec:criteria}). Both are reported.
\end{enumerate}
We leave out MESA's construct criteria (B01 to B04) and the extrapolation criteria B16 and B23, so Admissible here covers only the seven selected criteria, a narrower profile than MESA's mandatory one; the engine checks only that the transfer justification exists, names an author, and matches the target scope, not what it argues.

\subsection{Related Work}
Table~\ref{tab:paper2_related_work} compares related approaches. Assurance cases for machine learning \cite{hawkins_amlas_2021} and assurance contracts \cite{sljivo_assurance_contracts_2014} link requirements to evidence and require assumptions to be discharged before reuse. Work on evidence reuse in pipelines \cite{mukhopadhyay_nalla_evidence_reuse_2026} and on assurance after retraining \cite{picardi_transfer_assurance_2023} checks the integrity and validity of reused artifacts. VNVSpec \cite{arief_vnvspec_2026} links requirements to test runs and keeps inconclusive results when data are missing. Benchmark audits \cite{ding_spec_oracle_gaps_2026,openai_swebench_audit_2026}, BetterBench \cite{reuel_betterbench}, and Evaluation Cards \cite{ghosh_evaluation_cards_2026} assess benchmarks and their reporting, and model selection methods such as DAMSA--SMS \cite{singh_damsa_sms_2026} rank candidates on benchmark metrics. Automated transcript analysis \cite{mohl_transcript_flaws_2026} flags validity problems in agent runs, and the Agentic Benchmark Checklist \cite{zhu_agentic_benchmark_checklist_2025} lists review items for agentic benchmarks.

Automated audits \cite{tu_benchguard_2026,wang_auto_benchmark_audit_2026,wang_benchjack_2026} have agents find defects in tasks, tests, and environments. Table~\ref{tab:audit_map} maps their categories to MESA criteria; the mapping is one model's judgment from each paper's description. Audits supply checks for B05, B07, B28a, and B28b. The review also covers the construct, uncertainty reporting (B12), training exposure (B14a, B14b), and documentation, and turns the evidence into a decision on a requirement. MESA covers every shared concern by construction, so the mapping shows what audits leave out, not that MESA works. We treat audits and the review as complements, and compare an audit only decision with MESA's (Section~\ref{sec:evaluation}).

\begin{table}[t]
\caption{Audit categories and the MESA criteria they can supply checks for.}
\label{tab:audit_map}
\centering
\scriptsize
\setlength{\tabcolsep}{4pt}
\begin{tabularx}{\textwidth}{@{}>{\raggedright\arraybackslash}p{6.2cm}>{\raggedright\arraybackslash}X@{}}
\toprule
\textbf{Audit category} & \textbf{MESA criteria (fit)} \\
\midrule
BenchGuard ground truth, evaluation, instruction \cite{tu_benchguard_2026}; ABA instruction and evaluation \cite{wang_auto_benchmark_audit_2026} & B05 and B07 (full to partial); B06, B10, B27 in part \\
BenchGuard environment; ABA environment & B18, B28b (partial) \\
BenchJack isolation failure, answers shipped with the test, code execution in the evaluator, trusting untrusted output \cite{wang_benchjack_2026} & B28a, B28b, B07 (full) \\
BenchJack judge prompt injection, weak string matching, evaluation logic gaps & B07 (full) \\
BenchJack excessive permissions & B28b (partial) \\
No audit category & B01 to B04, B08, B09, B11, B12, B13, B14a, B14b, B15 to B17, B19 to B26, B29 \\
\bottomrule
\end{tabularx}
\end{table}

None of these turns a structured review of an external benchmark into a decision on a stated requirement with recorded reasons. The design question that MESA leaves open, and that this paper addresses, is which recorded information lets such a review count for a requirement, and how the decision and its explanation change when that information is missing or wrong.

\begin{table}[t]
\caption{Related assurance frameworks, specifications, and benchmark audits.}
\label{tab:paper2_related_work}
\centering
\scriptsize
\setlength{\tabcolsep}{3pt}
\begin{tabularx}{\textwidth}{@{}l>{\raggedright\arraybackslash}p{2.2cm}>{\raggedright\arraybackslash}p{2.3cm}>{\raggedright\arraybackslash}p{2.3cm}>{\raggedright\arraybackslash}X@{}}
\toprule
\textbf{Framework} & \textbf{Input Req. Format} & \textbf{Evidence Source} & \textbf{Missing Evidence} & \textbf{Relation to this study} \\
\midrule
AMLAS \cite{hawkins_amlas_2021} & ML safety requirements & Lifecycle assurance data & Not specified in cited work & Specifies assurance argumentation linking safety requirements to training and verification artifacts; does not define a procedure for qualifying external coding benchmarks. \\
Assurance Contracts \cite{sljivo_assurance_contracts_2014} & Assumption and guarantee contracts & Component safety claims & Assumptions must be discharged & Generates safety case argument fragments whose assumptions must be discharged for reuse; does not address confidence bounds or benchmark reviews. \\
Safety Case Impact \cite{lawford_safety_case_impact_2017} & Safety case arguments & System change diffs & Triggers reassurance & Assesses the impact of system modifications on existing safety cases rather than initial qualification of external benchmarks. \\
Transfer Assurance \cite{picardi_transfer_assurance_2023} & Target domain specs & Retrained model test sets & Reassesses affected artifacts & Retains parts of an assurance case and reassesses others after retraining; does not address evidence from external benchmarks. \\
VNVSpec \cite{arief_vnvspec_2026} & Machine readable DSL & Test execution \& verifiers & Marked as inconclusive & Links requirements to test runs and formal verification with inconclusive results; does not review external benchmarks. \\
Evidence Reuse \cite{mukhopadhyay_nalla_evidence_reuse_2026} & Target assurance goals & Pipeline traces \& DVC replays & Gate fails closed & States validity conditions and a release gate evaluated on 714 DVC pipelines; cannot detect undeclared dependencies and has no requirement thresholds for coding tasks. \\
Coding Agent Envelopes \cite{goswami_assurance_envelopes_2026} & Software change obligations & Typed inference graphs \& ASTs & Reports solver infeasibility & Derives supplied change obligations with forward chaining and a constraint solver at run time rather than judging benchmark evidence at design time. \\
NIST Guidance \cite{nist_ai_800_2_2026} & AI risk management goals & Benchmark leaderboards & Recommends documentation & Gives guidance for defining measurement targets; does not specify a decision rule linking benchmark scores to requirements. \\
Oracle Auditing \cite{ding_spec_oracle_gaps_2026} & Prompt docstrings & Hidden test suites & Flags implicit constraints & Analyzes overconstrained tests and specification gaps in coding benchmarks; does not link results to project requirements. \\
Model Selection in SE \cite{hou_llm4se_survey_2024} & SE task descriptions & Published benchmark scores & Disclosed as limitations & Surveys LLMs for software engineering; does not provide a decision rule for requirements acceptance. \\
DAMSA--SMS \cite{singh_damsa_sms_2026} & Deployment criteria & Candidate measurements & Eligibility screening & Applies Pareto screening, ranking, and sensitivity analysis; does not use benchmark reviews to withhold decisions. \\
ECBD \cite{liu_ecbd_2024} & Intended use specs & NLP evaluation rubrics & Flags unjustified choices & Uses intended use and validity warrants to guide benchmark authoring; does not judge requirements. \\
BetterBench \cite{reuel_betterbench} & Benchmark evaluation rubrics & Benchmark code \& metadata & Zero score on rubric criteria & Assesses 46 criteria on benchmark quality rather than decisions on candidate models. \\
Evaluation Cards \cite{ghosh_evaluation_cards_2026} & Standardized schemas & Benchmark \& model metadata & Displays omissions & Reports benchmark, run, and model metadata with comparability signals; has no decision thresholds. \\
Beyond the Hype \cite{pizarro_beyond_the_hype_2026} & Organizational adoption goals & Empirical evaluation metrics & Flags adoption risks & Adapts GQM for organizational adoption rather than requirement decisions from benchmark reviews. \\
\midrule
\textbf{This work} & \textbf{Requirement with threshold and evidence conditions} & \textbf{Benchmark observations and MESA rating forms} & \textbf{Inquiry Required; no decision} & \textbf{Applies MESA's traceability rules as records and checks; reports input integrity, evidence status, and decision separately with reasons.} \\
\bottomrule
\end{tabularx}
\end{table}


\section{Method}
\label{sec:method}

\subsection{Case and Rating Form Records}
A \emph{case} is one requirement, one coding agent (a model in its tool, a configuration in the records), and the observations used to judge it. Table~\ref{tab:records} lists what a case record contains. Each case record carries a schema version; the engine refuses unknown versions.

\begin{table}[t]
\caption{Contents of a case record and the level at which each part is used.}
\label{tab:records}
\centering
\scriptsize
\setlength{\tabcolsep}{3pt}
\begin{tabularx}{\textwidth}{@{}>{\raggedright\arraybackslash}p{2.5cm}>{\raggedright\arraybackslash}X>{\raggedright\arraybackslash}p{2.6cm}@{}}
\toprule
\textbf{Part} & \textbf{Content} & \textbf{Used at} \\
\midrule
Requirement & Identifier, target repository, task set, configuration, threshold $\theta$, confidence level, minimum number of observations $n_{\min}$ & Input integrity \\
Evidence conditions & For each criterion: minimum level, required check outcome, accepted scope levels, required fields; whether execution provenance is required & Input integrity (form); evidence status (use) \\
Configuration and tasks & Configuration identifier; declared task identifiers & Input integrity \\
Observations & Task, configuration, repository, resolved or not, patch (task, path, hash), other artifact hashes, execution provenance (session, log, log hash) & Input integrity (identity, hashes); evidence status (provenance present) \\
Rating forms & Criterion (B05, B07, B12, B14a, B14b, B28a, B28b) and concern (C05 \dots C28), definition path and hash, applicability and reason, level, outcome, check kind, scope (release, configuration, task, or subset, with covered configurations and tasks), check procedure, evidence location, detection limits, task defects, reasons, who rated and who signed & Input integrity (identity, definition); evidence status (conditions) \\
Transfer justification & Text, author, source scope, target scope & Requirement decision \\
\bottomrule
\end{tabularx}
\end{table}

\subsection{Three Decision Levels}
\label{sec:levels}
The engine evaluates a case in three levels and keeps their outputs apart.

\textbf{Input integrity} (valid or rejected). The engine rejects a case before judging capability when an observation belongs to another configuration, task, or repository; when a declared task has no observation; when a patch, artifact, or log no longer matches its recorded hash; when an observation cites another task's patch, or two observations cite the same patch; when a rating form cites the wrong concern or an unfrozen definition; when a criterion has two records; when the policy lists no evidence conditions; or, for the study, when the policy differs from the pinned study policy. A rejected case receives no evidence status and no decision.

\textbf{Evidence status} (Admissible, Inquiry Required, Inadmissible, or not applicable). For each criterion $c$ named by requirement $r$, the status $v(c,r)$ is
\begin{itemize}[leftmargin=*]
    \item \emph{not applicable} if the rating form marks $c$ as Not applicable and gives a reason;
    \item \emph{Inadmissible} if the level is \emph{Contradicted} or the check failed, the defect has a cited location, and any recorded counterexample was executed;
    \item \emph{Admissible} if the level meets the required minimum, the recorded check passed, the scope level is accepted, the scope covers the configuration and every task, the required fields (procedure, check kind, evidence location, detection limits) are present, and the roles of the author and the model are recorded;
    \item \emph{Inquiry Required} otherwise, including a missing record, a defect without a location, a Not applicable mark without a reason, and a Not applicable mark on B05, B12, B14a, or B14b, which apply always (B14b whenever a configuration is named).
\end{itemize}
The aggregate status is Inadmissible if any criterion is Inadmissible; otherwise Inquiry Required if any criterion requires inquiry, required execution provenance is missing, or fewer than $n_{\min}$ tasks remain after defective tasks are removed; otherwise not applicable if every criterion is not applicable; otherwise Admissible. A defect therefore takes precedence, but the reasons from every blocking criterion are returned. MESA levels remain labels; nothing is scored or averaged.

\textbf{Decision} (Accept, Reject, or Inconclusive, or no decision). Tasks that a rating form lists as defective, each with a cited location, are removed before any score is seen. For $s$ resolved tasks out of $n$ remaining observations, the 95\% Wilson score interval \cite{wilson_probable_inference_1927,brown_cai_dasgupta_2001} is
\begin{equation}
\label{eq:wilson}
\text{LB}, \text{UB} = \frac{\hat{p} + \frac{z^2}{2n} \pm z \sqrt{\frac{\hat{p}(1-\hat{p})}{n} + \frac{z^2}{4n^2}}}{1 + \frac{z^2}{n}}, \qquad \hat{p} = s/n,\; z = 1.95996.
\end{equation}
The score decision $D_{\text{perf}}$ is Accept if $\text{LB} \ge \theta$, Reject if $\text{UB} < \theta$, and Inconclusive otherwise. Then
\begin{equation}
\label{eq:decisions}
D_{\text{src}} = \begin{cases} D_{\text{perf}} & \text{if the evidence is Admissible} \\ \text{no decision} & \text{otherwise}, \end{cases}
\qquad
D_{\text{req}} = \begin{cases} D_{\text{src}} & \text{if transfer is justified} \\ \text{no decision} & \text{otherwise.} \end{cases}
\end{equation}
Transfer is justified when the case carries a justification with an author, its source scope matches the source observations, and its target scope matches the requirement's repository and the configuration. Inadmissible evidence never yields a Reject: defects can lower measured success as well as raise it, for example tests that reject a correct fix, so a low resolve rate on defective evidence does not show that the coding agent fails the requirement. Missing information never counts as failure. The output lists the reasons from all three levels, including the interval when it spans the threshold, and reports the score with and without the removed tasks.

\subsection{Selected Criteria and Study Policy}
\label{sec:criteria}
The evaluation uses seven MESA criteria (Table~\ref{tab:criteria}). They decide whether a benchmark resolve rate is a sound measurement for this requirement, and they follow the scoring, generalization, and extrapolation steps of Kane's validity argument \cite{kane_validation}. We did not choose B01 to B04, because the requirement defines the ability itself; B16 and B23, which concern extrapolation, because the transfer justification stands in for them and RQ2 tests that step directly; or the rest, which concern reporting, licensing, maintenance, fairness, or societal effects. The selection is this study's own choice; another requirement would call for other criteria, and each criterion added makes the strict policy harder to meet.

\begin{table}[t]
\caption{Selected criteria and what each asks.}
\label{tab:criteria}
\centering
\scriptsize
\setlength{\tabcolsep}{4pt}
\begin{tabularx}{\textwidth}{@{}lX@{}}
\toprule
\textbf{Criterion} & \textbf{What it asks} \\
\midrule
B05 & Do any tests reject a correct solution? \\
B07 & Do any tests or scoring code accept a wrong solution? \\
B12 & Does the benchmark report its own uncertainty? \\
B14a & Does the benchmark guard against training exposure of its tasks? \\
B14b & Could one named model have seen the tasks in training? \\
B28a & Is the grading environment controlled, so that the agent cannot reach answers or change the scoring? \\
B28b & Is the environment the benchmark ships for agents controlled and isolated from answers? \\
\bottomrule
\end{tabularx}
\end{table}

\emph{Training exposure.} A model may have read a task's issue and its fix during training and remember the fix instead of working it out. B14a asks what the benchmark does against this. B14b asks it for one named model; a task created after the model's stated training cutoff counts as a held out check. The source tasks are old and public, so exposure cannot be ruled out; the target tasks were fixed after the stated training cutoffs of both models (Section~\ref{sec:evaluation}). \emph{Environments.} B28a concerns the grading environment, which runs the hidden tests on a patch. B28b concerns the agent environment that the benchmark ships, in which an agent works on a task. Neither is the environment of the sessions run here, which enters through the validity rules.

Two policies state which levels count as enough. The \emph{strict} policy (the study policy) requires level \emph{Verified}, a passed check, a scope that covers the coding agent and tasks, and a recorded procedure, check kind, evidence location, and detection limits, for all seven criteria. The check outcome is the result of the check that the procedure describes, whether the benchmark sources report it or the study runs it; the check kind says which. The policy is pinned by its digest, and the study entry point rejects any case with another policy. A \emph{Verified} level without a passed outcome, with a scope that omits a coding agent or task, or without detection limits therefore leads to Inquiry Required. The \emph{lenient} policy requires B05, B28a, and B28b at \emph{Verified} and B07 at \emph{Documented}, each with a passed check, and leaves out B12, B14a, and B14b. It shows the decision when training exposure and the benchmark's own statistics are ignored; its results are reported apart from the strict policy.

\subsection{Comparison Methods}
For each coding agent we compare decisions on the 20 source tasks, each with the interval of Eq.~\ref{eq:wilson} and the threshold. Method~1 (M1) applies Eq.~\ref{eq:wilson} to all evaluated source observations. Method~2 (M2) applies it to the observations that match the requirement's repository; it is the resolve rate alone, with no review. In the selected scope every source observation is a SymPy task, so M1 and M2 coincide. \emph{Audit only} leaves out the tasks the executed checks show to be broken, with no criteria gate, as automated audits use their findings; it equals MESA's decision before the evidence status. \emph{ABA as a second audit} leaves out the tasks that ABA's published findings flag with a major finding, which are two of the 20 source tasks (\texttt{sympy\_\_sympy-13031} and \texttt{-18199}), and decides on the rest. Method~3 (M3) is MESA: it uses M2's observations less the tasks whose rating forms show broken tests, adds the evidence status, and gives the decision only when the status is Admissible, under the strict and the lenient policy. For each blocking criterion it reports what its form says is missing, and it names each task left out with the reason. The fifth method is ABC \cite{zhu_agentic_benchmark_checklist_2025}, as its paper describes it. Its 33 items that MESA's concern mapping assigns to the concerns of the seven criteria are answered as met, not met, not applicable, or unknown, each with a quote and a location. Each item scores 1 when met and 0 otherwise, and the score of a part is the average over its applicable items (its Section~3). The decision comes from the resolve rate on all 20 source tasks with the same interval, and each unmet or unknown item is a caveat (its Appendix F, items R.10 and R.11). ABC never withholds a decision, so its decision equals M2's apart from its caveats.

\subsection{Explanations}
Each output names its level and its cause. For example, a constructed case with 20 resolved SymPy observations, all seven criteria rated \emph{Verified}, and a B07 record without a recorded check outcome yields: input integrity valid; evidence status Inquiry Required, with B07 listed as the blocking criterion; and no decision on the source tasks or on the requirement. Removing the execution provenance from one observation adds \emph{execution provenance missing} for that task, while deleting a log file or editing a patch after collection rejects the case for a missing or stale artifact.

\section{Prepared Evaluation}
\label{sec:evaluation}
This section describes the evaluation protocol. No study session, no replay, and no new simulation has been run. The eligibility checks, the checks of the published environments, sandbox checks of both agent tools, and a rehearsal of the session procedure for both coding agents on one task outside the selection have been run. The exploit attempts for B07, the grading check for B28a, and the ABC answers were made before the forms are signed; the two rating runs precede signing as well, and the author's signatures are pending.

\subsection{Scope and Session Budget}
The evaluation uses one repository, SymPy, and two coding agents. Each is a model in its tool, called a configuration in the records: Gemini 3.8 Flash Medium in the Antigravity command line tool (\texttt{CONFIG\_01\_GEMINI\_FLASH\_MED}) and GPT-6.1 Sol at medium reasoning effort in the Codex command line tool (\texttt{CONFIG\_02\_SOL\_MED}). The threshold is $\theta = 0.70$. A \emph{session} is one coding agent working on one task for at most 15 minutes, with no internet and no human input. Table~\ref{tab:sessions} gives the session budget. Each phase has its own slots in a session ledger, and a slot is reserved by exclusive file creation before launch, so concurrent launches cannot exceed a phase limit. Failed starts and interruptions keep their slots and count. A reserve slot requires a recorded infrastructure failure, each failure can be replaced once, retries are never automatic, and there is no 91st slot. Each session starts from a script written when its slot is opened. The script runs the tool once in a prepared workspace with fixed flags, gives it no input after the prompt, and stops it at 900 seconds. The author launches it, or Claude Opus 5.5 at the author's request; the tooling itself never starts a model. A session opens only after the protocol, rating forms, ABC record, policy, prompt, requirement, and task selection are frozen and the author has authorized its phase in a record that names the frozen manifest's hash. The pilot phase runs the procedure on one task per coding agent after the freeze; it is not used for RQ1 or RQ2, and its transcripts count in the session validity control.

\begin{table}[t]
\caption{Session budget of the prepared evaluation.}
\label{tab:sessions}
\centering
\scriptsize
\begin{tabular}{lr}
\toprule
\textbf{Phase} & \textbf{Sessions} \\
\midrule
Pilot: one task $\times$ two coding agents & 2 \\
Source: 20 tasks $\times$ two coding agents & 40 \\
Target: 20 tasks $\times$ two coding agents & 40 \\
Reserve for infrastructure failures & up to 8 \\
\midrule
Maximum & 90 \\
\bottomrule
\end{tabular}
\end{table}

\subsection{Requirement, Source Tasks, and Target Tasks}
The requirement is the one in Section~\ref{sec:intro}. It names SymPy as its repository and also states its task set (the 20 source tasks, less any with shown broken tests), the coding agent, the confidence level (95\%), and the minimum number of observations $n_{\min} = 9$. Under the 95\% Wilson rule at $\theta = 0.70$, nine resolved of nine give a lower bound of 0.701 and eight of eight give 0.676, so nine is the smallest number of tasks at which Accept is possible. Claude Opus 5.5 drafted the requirement and one transfer justification per coding agent, and the author signs them.

An \emph{issue} is a bug report on SymPy's GitHub tracker that the maintainers later fixed, adding tests with the fix. A \emph{task} is one issue prepared for a coding agent: the issue text and SymPy's code as it was before the fix, with the tests from the real fix hidden from the agent. For example, the issue \texttt{sympy\_\_sympy-21930} reports that LaTeX printing of the square of \texttt{Bd(a)} gives an invalid double superscript. We use it only in a rehearsal, outside the selection. The agent gets that text and the code from before the fix; the maintainers' tests decide whether it is resolved.

\emph{Source tasks} are the evidence a practitioner would read: the 75 SymPy tasks of SWE-bench Verified, saved in full from a fixed revision of the dataset, with issues from 2016 to 2023. They are old and public, so the models may have seen them in training. \emph{Target tasks} are the direct test: SymPy pull requests merged from 1 May to 5 November 2026 that close a linked issue and change both a test file and a source file. The window starts after the later of the two training cutoffs stated by the providers (March 2026 for Gemini 3.8 Flash; 30 April 2026 for GPT-6.1 Sol). The search was fixed before it was run; of 322 merged pull requests, 108 meet the rule. Candidates are checked in a fixed order: for each, the GitHub interface confirms the issue, the merge and its date, and the base commit; the pull request is split into the reference patch and its tests, and the FAIL\_TO\_PASS and PASS\_TO\_PASS tests are found by running the tests before and after the reference patch. The pull requests that fix the target issues were merged from 13 June to 5 November 2026, after both models' stated training cutoffs. If the stated cutoffs hold, the fixes were not in the models' training data; five target issues (27443, 13365, 21772, 19324, and 26652) were opened before the cutoffs, so the models may have seen their text. The cutoffs are the providers' claims, not checks.

A task is eligible when, with the source unchanged, every FAIL\_TO\_PASS test fails and every PASS\_TO\_PASS test passes, and with the reference patch every test passes. A run that produces no report is a harness failure, kept as a record, and leaves the task pending. All 75 source tasks are eligible. The first 20 target candidates in the order are all eligible, so no candidate was skipped. Source tasks run in the published SWE-bench environments; target tasks run in environments built for this study with pytest, which current SymPy uses. A session's patch counts as resolved when every FAIL\_TO\_PASS and PASS\_TO\_PASS test passes in this harness, the session stayed within 900 seconds, and the session is valid (Section~\ref{sec:sessions}). A session over a limit is valid but not resolved.

Tasks are ordered by a hash of a fixed label and the task identifier. The pilot task and the 20 source tasks are the first 21 eligible tasks of the source pool; the 20 target tasks are the first 20 eligible tasks of the target pool. Model results are never consulted. The rehearsal task, \texttt{sympy\_\_sympy-21930}, is the last of the 75 source tasks in this order and outside the selection. Pilot, source, and target sets must not share a task or an issue or pull request number; the selection step checks this. The check is partial, because source tasks record only their pull request number, not the issues they close. A shortfall would block selection; no repository or task is substituted.

\subsection{Rating Forms, Exploit Attempts, ABC, and Signing}
Eight rating forms cover B05 for the pilot and the 20 source tasks, B07 for the 20 source tasks, B12, B14a, B28a, and B28b for the benchmark release, and B14b for each coding agent, because exposure of a model depends on the model. B07, B28a, and B28b concern the benchmark release, that is, its grading harness and published environments, not the sessions run here. A form records applicability, evidence locations with short quotes, the check performed, its kind and outcome, scope, detection limits, task defects, reasons, what was missing, and the roles of the author and the model. A form with a level but without evidence and reasons is rejected.

Claude Sonnet 5.5 writes the eight forms in two separate runs, A and B, from the same fixed prompt, saved before the runs. The prompt gives it the rubric, the saved sources, and the check records, including the exploit attempts, and excludes the decision code, the policies, and other ratings. Both runs are kept unchanged. The author reviews and signs the forms of run A; run B is the repeat, and the agreement of the two runs is recorded per form for level, outcome, check kind, and task defects. A third Claude Sonnet 5.5 run answers the ABC items from a fixed prompt with the same inputs, and the author reviews and signs that record too. The author's changes to the forms and the ABC record at signing are listed field by field and reported as a description, not as a finding. The freeze refuses an unsigned form, an unsigned ABC record, or an unsigned requirement.

\emph{Exploit attempts for B07.} B07 asks whether the scoring can be gamed. In SWE-bench the scoring is each task's hidden tests, so it is gamed when a patch passes them without fixing the issue. For each of the 20 source tasks, Claude Sonnet 5.5 makes one attempt, from the issue text and the real fix only, never from the hidden tests or any session result. It lists the behaviors the issue asks for, writes an independent check of them, and writes at most three cheating patches, each meant to pass the hidden tests while breaking one listed behavior (for example, the real fix without one part), or states why none could be written. The tool runs the check on the code before the fix, where it must fail, and with the real fix, where it must pass; otherwise the check is invalid and no patch counts. It then grades each patch on the hidden tests and runs the check on it. A defect is shown on a task when the check is valid, the hidden tests pass the patch, and the check fails on it. MESA counts it as a defect of the task's tests only if a behavior the patch breaks is stated in the issue text, because a check may also test behaviors inferred from the real fix; the rating runs judge this from the issue text, quoting it. Every attempt is recorded, whatever the outcome. The B07 form lists each shown defect, and for the other tasks the attempts are its check, with the limits that they are one model's construction, one attempt per task and at most three patches, so that a failure to find a defect does not show that none exists. The attempts were made before the rating forms were written. A cheating patch passed the hidden tests and failed the check on 17 of the 20 tasks. Both rating runs count six of them as defects of the tests, because the patch breaks a behavior the issue states (\texttt{sympy\_\_sympy-13878}, \texttt{-15809}, \texttt{-18199}, \texttt{-21847}, \texttt{-22914}, and \texttt{-23824}), and both set aside the same 11, whose patches break only behaviors inferred from the real fix.

\emph{Grading check for B28a.} B28a asks whether the agent can change the scoring. In SymPy's grading the patched repository runs its own test script, so a patch can print passing lines. Claude Sonnet 5.5 wrote three patches for each of three source tasks, one per test runner layout (\texttt{sympy\_\_sympy-13031}, \texttt{-20590}, and \texttt{-24539}), that change only the test running files (\texttt{bin/test}, \texttt{runtests.py}) and no source code. The SWE-bench harness graded all nine as resolved. B28a is therefore \emph{Verified} with outcome failed. The defect lies in the grading of the whole benchmark, so the form lists no task. The limits are three tasks, three patches each, and one model.

\emph{Tasks with broken tests.} A form may list single tasks whose tests are shown to be wrong, with a cited location and a quote: under B05 when the tests reject a correct solution, and under B07 when they accept a wrong one. A cheating patch counts only when it was executed. M3 leaves such a task out before any score is seen, names it with the reason, and reports the score with and without it; M1, M2, and ABC keep every task. A defect of the whole benchmark, such as the grading defect, blocks M3. With the rating forms as written, before the author's signature, M3 leaves out six of the 20 source tasks: \texttt{sympy\_\_sympy-18199} under B05 and B07, and \texttt{-13878}, \texttt{-15809}, \texttt{-21847}, \texttt{-22914}, and \texttt{-23824} under B07, so 14 remain.

\emph{B05 and the selected tasks.} For B05, the 2024 annotation of SWE-bench Verified screened every task for tests that reject valid solutions \cite{openai_swebench_verified_2024}. A 2026 audit names \texttt{sympy\_\_sympy-18199} as having tests wider than its issue \cite{openai_swebench_verified_contamination_2025}, and that task is one of the 20 source tasks. Claude Opus 5.5 graded a patch limited to the reference lines for the stated issue: the hidden tests rejected it, with the same result as the unchanged source, while the reference patch passed. The audit examined 138 SWE-bench Verified problems, all of which had passed the 2024 screening, and reports flawed tests in at least 59.4\% of them; the 2024 screening therefore missed defects that a later check found, which is a detection limit of the screening. The rating run rates B05 \emph{Verified} for the pilot and the 20 source tasks, with \texttt{sympy\_\_sympy-18199} as its one task defect, and records the outcome as passed: the outcome describes the tasks in scope after the listed defect is left out. The audit's rate of flawed tests names no task here, so it is a detection limit.

\emph{ABA on the same tasks.} ABA's published findings flag \texttt{sympy\_\_sympy-13031} and \texttt{-18199} with a major finding. The forms list six defects. Only \texttt{-18199} is in both, with the same finding as B05: the tests require composite modulus support that the issue does not ask for. The other five come from executed B07 exploits that ABA's audit did not find, and ABA's finding on \texttt{-13031} is not in the forms.

\emph{ABC as answered.} As answered by the model run, ABC has 8 items met, 9 not met, 10 not applicable, and 6 unknown. The scores of its parts are 5 of 8 for task validity (T), 2 of 11 for outcome validity (O), and 1 of 4 for reporting (R). The 15 items not met or unknown are its caveats; the decision comes from the resolve rate on all 20 source tasks. The author's changes at signing are reported.

\subsection{Sessions, Validity, and Isolation Checks}
\label{sec:sessions}
Before a session, the sandbox check runs each tool with the session flags in a scratch folder and records whether network access and reads outside the folder are blocked; network isolation is reported as enforced only where blocking is shown. A rehearsal runs both coding agents once on the task outside the selection through the same open, run, close, and grading steps, with its own ledger and records; its results are not study results and enter no decision, and a problem it shows is fixed before the freeze. The session then runs from its script, and \texttt{close} reads the logs, refuses a log that starts before its slot was reserved, grades the patch on the hidden tests in the harness, and archives the workspace.

A session is \emph{invalid} when its commands or file paths name study material (the study records and folder, other workspaces, the repository history with later fixes, or the agent tools' own data), its tool outputs show oracle content such as test lists, it runs a network command, it calls a tool outside its tool's local set (for example web search, browser, or subagent tools), it receives a message after the prompt, or its patch changes a test running file (\texttt{bin/test}, any \texttt{runtests.py} or \texttt{conftest.py}, \texttt{pytest.ini}, \texttt{setup.cfg}, \texttt{tox.ini}, or SymPy's own \texttt{pytest.py}). No reference fix of the selected tasks touches these files. Study material is matched by marker lists in the session code. An invalid session keeps its slot, is recorded with the matched evidence, and counts as not resolved. These rules are a control, not a research question. Agents working on SWE-bench tasks have retrieved fixes from git history or the web \cite{jain_cursor_reward_hacking_2026,stein_debugml_cheating_2026}, and automated scanners flag such access in transcripts \cite{mohl_transcript_flaws_2026}; the last rule follows from the grading check. The transcripts of all sessions are counted per coding agent by kind of invalid session, and the resolve rate and decision are reported with invalid sessions counted as not resolved and by test result alone.

\subsection{Replay as a Test of the Implementation}
The replay tests the implementation on collected records; it answers no research question. For each source task of one coding agent, it evaluates the case as collected, then two changed copies (Table~\ref{tab:replay}). Each copy records its parent's digest and lists its changes, and is written once. Each changed copy must change what it targets: the provenance reason for that task, or a rejected input. Source tasks under three conditions are repeated observations on the same tasks, not independent tasks.

\begin{table}[t]
\caption{Replay conditions and the level each is designed to change.}
\label{tab:replay}
\centering
\scriptsize
\setlength{\tabcolsep}{4pt}
\begin{tabularx}{\textwidth}{@{}l>{\raggedright\arraybackslash}X>{\raggedright\arraybackslash}p{3.8cm}@{}}
\toprule
\textbf{Condition} & \textbf{Change to the case record} & \textbf{Level designed to change} \\
\midrule
Original evidence & None & None \\
Provenance removed & The task's execution provenance removed (for example, the log reference) & Evidence status: Inquiry Required \\
Identity mismatch & The task's identifier changed, or its observation cites another task's patch & Input integrity: rejected, no decision \\
\bottomrule
\end{tabularx}
\end{table}

The designed changes follow from the rules in Section~\ref{sec:levels} and are checked on software fixtures; they are not findings. The replay asks whether the checks detect the changes in real artifacts and whether the reasons name the changed task and element. Under the strict policy the aggregate evidence status may already be Inquiry Required, so a condition can show its effect only on the criterion or reason it targets; the replay rows record both.

\subsection{Readiness Controls}
A readiness report lists the implemented checks, the fixture tests, and the capabilities that still need experimental verification. Checking file locations, telemetry fields, or container flags is reported as configuration validation. It does not establish observed model identity, the 900 second limit, isolation from reference fixes and tests, network isolation, or resource limits; each of these needs an observation from a recorded session.

\subsection{Planned Analysis}
\label{sec:planned}
For RQ1, we will report, for each coding agent, the decision of each method on the source tasks: the resolve rate alone (M2), audit only, ABA, MESA under the strict and the lenient policy, and ABC. For MESA we report the evidence status, each blocking criterion with what its form says is missing, and the tasks left out; for ABC, the part scores and the caveats. We will report the inputs: for each criterion, its level, outcome, and check kind, and whether a \emph{Verified} level rests on a check the sources report or one the study ran; the exploit attempts and how many the rating runs count as defects; the grading check; the B14b ratings for the two named models; whether run B repeats the scope judgments, the task defects, and the B14b ratings of run A; and what stayed missing. For RQ2, we will report, for each coding agent and each decision from the source tasks, the direct test's decision and whether the two agree; agree means the same decision, fixed in advance. For MESA we also report the decision it would give on its kept tasks when it withholds one. We will report the source minus target resolve rate with its 95\% interval (Newcombe's method), and the direct test without the five target tasks whose issues were opened before the cutoffs. Two coding agents give two cases, not a rate, and the target tasks differ from the source tasks in more than exposure, so a difference cannot be split into training exposure and task difficulty. The control reports, for each coding agent, the invalid sessions by kind in the pilot, source, and target phases, and the resolve rate and decision with invalid sessions counted as not resolved and by test result alone; whether invalid sessions occur is not known in advance. The replay is a test of the implementation, not a research question: a changed copy that does not change the reason it targets on collected records counts against the implementation.

The following is known in advance. With 20 tasks at 0.70 under the 95\% Wilson rule, Accept needs 19 or 20 resolved and Reject needs 9 or fewer, so the resolve rate alone decides only when the observed rate is far from 70\%. If tasks are removed, these cut points move, and below $n_{\min} = 9$ remaining tasks the evidence status is Inquiry Required. With 18 tasks (ABA), Accept needs 17 or more resolved and Reject 8 or fewer. With the six tasks left out, MESA's decision before the evidence status and the audit only decision use 14 tasks: Accept needs 14 of 14 resolved, and Reject 6 or fewer. ABC decides on all 20 tasks, with 15 caveats (8 items met, 9 not met, 10 not applicable, 6 unknown).
With the forms as written before the author's signature, the strict policy gives Inadmissible, because B14a and B28a have outcome failed, with inquiry on B12 and B14b for both coding agents; the lenient policy gives Inadmissible, because of B28a. MESA therefore withholds its decision under both policies for both coding agents, while the other methods decide. The six tasks left out are \texttt{sympy\_\_sympy-13878}, \texttt{-15809}, \texttt{-18199}, \texttt{-21847}, \texttt{-22914}, and \texttt{-23824}. The author's review at signing may change these levels, and the decisions are then fixed before any session. A decision on the source tasks does not become a decision on the requirement without a signed transfer justification.

\section{Completed Calculations and Simulation}
\label{sec:completed}
This section reports calculations completed before the sessions. They use the interval rule of Eq.~\ref{eq:wilson}, with decisions named Accept, Reject, and Inconclusive.

\subsection{Sample Size}
Table~\ref{tab:power_analysis} gives exact binomial probabilities of each decision under Eq.~\ref{eq:wilson} at $\theta = 0.70$. With 30 observations and a true success rate of 0.75, an Accept decision occurs with probability 9.8\% and an Inconclusive decision with 89.4\%. At a true rate of 0.85, the probability of Accept first exceeds 80\% at $N = 60$ (81.9\%); because of the discrete binomial distribution it falls below 80\% again at some larger sizes, such as 78.8\% at $N = 62$. These probabilities assume admissible evidence; review findings can only add withheld decisions.

\begin{table}[t]
\caption{Exact binomial decision probabilities under 95\% Wilson bounds ($\theta = 0.70$).}
\label{tab:power_analysis}
\centering
\scriptsize
\setlength{\tabcolsep}{6pt}
\begin{tabular}{rcccccc}
\toprule
\textbf{Sample Size ($N$)} & \multicolumn{2}{c}{\textbf{True $p = 0.75$}} & \multicolumn{2}{c}{\textbf{True $p = 0.65$}} & \multicolumn{2}{c}{\textbf{True $p = 0.85$}} \\
\cmidrule(lr){2-3} \cmidrule(lr){4-5} \cmidrule(lr){6-7}
& \textbf{Accept} & \textbf{Inconclusive} & \textbf{Reject} & \textbf{Inconclusive} & \textbf{Accept} & \textbf{Inconclusive} \\
\midrule
$N = 15$ & 8.0\% & 90.3\% & 11.3\% & 87.3\% & 31.9\% & 68.1\% \\
$N = 30$ & 9.8\% & 89.4\% & 12.6\% & 86.6\% & 52.4\% & 47.5\% \\
$N = 45$ & 9.4\% & 90.2\% & 12.1\% & 87.5\% & 63.9\% & 36.1\% \\
$N = 60$ & 14.8\% & 84.9\% & 17.1\% & 82.5\% & 81.9\% & 18.1\% \\
$N = 90$ & 16.5\% & 83.4\% & 18.8\% & 81.1\% & 92.6\% & 7.4\% \\
$N = 120$ & 23.3\% & 76.6\% & 25.0\% & 74.9\% & 98.1\% & 1.9\% \\
\bottomrule
\end{tabular}
\end{table}

\subsection{Assumed Review Accuracy}
A simulation (10,000 trials per scenario, $N = 30$, $\theta = 0.70$) varied the rate $d$ of defective benchmark evidence and the shift $\Delta$ between source and target success rates. The source success rate was 0.72 and rose to 0.86 when the evidence was defective. The review was \emph{assumed} to detect 85\% of defects, to raise false alarms on 5\% of sound evidence, and to lack documentation in 10\% of cases. These are model inputs, not measured accuracy of MESA reviews. M3 was compared with two controls that withhold the same number of M2 decisions: at random, and closest to the threshold \cite{el_yaniv_selective_classification_2010}.

\begin{table}[t]
\caption{False acceptance rate per case under assumed review accuracy ($N = 30$, $\theta = 0.70$, 10,000 trials).}
\label{tab:simulation_far}
\centering
\scriptsize
\setlength{\tabcolsep}{4pt}
\begin{tabular}{lccccccc}
\toprule
\textbf{Defect rate ($d$)} & \textbf{Shift ($\Delta$)} & \textbf{M2 cov.} & \textbf{M3 cov.} & \textbf{M2} & \textbf{Random} & \textbf{Proximity} & \textbf{M3} \\
\midrule
$0.10$ & $-0.15$ & 12.2\% & 6.4\% & 10.2\% & 5.4\% & 5.6\% & 4.7\% \\
$0.10$ & $-0.05$ & 12.9\% & 7.0\% & 10.6\% & 5.7\% & 6.3\% & 5.0\% \\
$0.20$ & $-0.15$ & 18.7\% & 6.7\% & 16.9\% & 6.1\% & 6.5\% & 5.2\% \\
$0.20$ & $-0.05$ & 17.1\% & 6.2\% & 15.5\% & 5.6\% & 6.0\% & 4.8\% \\
$0.30$ & $-0.15$ & 22.4\% & 6.8\% & 20.8\% & 6.4\% & 6.5\% & 5.4\% \\
$0.30$ & $-0.05$ & 22.2\% & 6.5\% & 20.5\% & 6.0\% & 6.3\% & 5.1\% \\
\bottomrule
\end{tabular}
\end{table}

At $d = 0.30$ and $\Delta = -0.15$, false acceptance fell from 20.8\% under M2 to 5.4\% under M3, against 6.4\% for random and 6.5\% for proximity withholding at the same 6.8\% coverage (Table~\ref{tab:simulation_far}). Most of the reduction comes from withholding decisions at all; the part due to the assumed review is about one percentage point. Among the decisions that remained, 79.2\% were still false acceptances, because the simulated review inspects benchmark defects and not the shift to the target. When the target was compliant ($d = 0.30$, $\Delta = 0$), M3 resolved 7.4\% of cases and wrongly rejected 1.42\%, against 0.44\% for random withholding. Lower false acceptance in noncompliant scenarios thus came with more abstention and false rejection in compliant ones; without a loss model, this is a trade between errors, not an overall improvement.

A sensitivity sweep of the same simulation over 576 combinations of assumed detection rate, false alarm rate, missing documentation rate, defect rate, and shift (2,000 trials each) shows that the result depends on the assumed accuracy. Without defects, the difference in false acceptance between M3 and random withholding ranged from $-0.40$ to $+0.45$ percentage points across the 72 noncompliant cells; with 30\% defects, M3 was lower in all 72 cells, by 0.20 to 1.50 points. Repeating the headline scenario over 100 seeds with a second implementation gave mean reductions of 0.88 points against random withholding (standard deviation 0.09) and 1.14 points against proximity withholding (standard deviation 0.10); these describe Monte Carlo variation under the assumed model, not uncertainty about real review accuracy.

\section{Threats to Validity and Limitations}
\label{sec:threats}

\textbf{Status of the evaluation.} No study session, no signed rating form or ABC record, and no replay has been conducted; eligibility, the checks on the published images, sandbox checks of both agent tools, and a rehearsal on one task outside the selection, whose results are not study results, have been run. The implementation is checked by software tests on constructed fixtures; these show that the code follows its rules, not that the rules help requirement decisions. Runtime model identity, the time limit, and isolation remain unverified.

\textbf{Sessions on the host.} Agents run on the host without CPU or memory limits. Reads outside the workspace are detected from the commands, file paths, and tool outputs in the logs, not prevented; a read that names no marked place and prints no oracle content is not detected. The two coding agents are isolated differently. In sandbox checks the Codex sandbox blocked network access from shell commands but allowed reads outside the folder, and its web search still returned results, first through its web tool and then through its apps server, so both are turned off; a later check found no web search tool. The Antigravity command line tool with all tools approved blocked neither. It runs every shell command in its sandbox without approval, as a custom agent with only shell and file tools; in a check of this setting, its sandbox blocked network access and reads and writes outside the folder, and the agent had no web search tool. Codex still offers subagent and other tools outside the shell, and any call to them makes a session invalid. Run without a prompt, the Antigravity tool cannot ask for approval, so a hook refuses a file tool access outside the workspace and a request to leave its sandbox, with a reason, and the run goes on, as Codex goes on after a refused command. Both tools are locked against updates, because an update after the freeze would change the tool. The reference fixes are on the same computer, in the dataset cache, the repository mirror, and the study records; Codex can read them, and only the markers detect such a read. Both agents find the same Python first: for source tasks Python 3.9.20 with mpmath 1.3.0, as in the images that grade them, and for target tasks Python 3.12.14 with the packages of the target images, which use Python 3.12.16; each also has pytest. Codex can still run other Python installations on the computer. The same invalidation rules apply to both coding agents, but a difference between their outcomes may partly reflect their isolation. The 900 second limit is enforced by the script, and a session over it does not count as resolved. Steps are counted from the logs and reported, not limited, because the two tools count them differently. The model and effort are those reported by each tool's own log, which is not attestation by the provider. The Antigravity tool keeps knowledge across conversations; it is moved aside before and after each session. Source tasks are graded in the published amd64 images under emulation with SymPy's runner, and target tasks in images built here with pytest; a difference between source and target outcomes may partly reflect this. Target tasks were not screened by people, unlike SWE-bench Verified tasks, and five target issues were reported before the training cutoffs; only their fixes are later.

\textbf{Internal validity.} The simulation assumes accuracy of the rating forms, independent missing documentation, defects that only inflate scores, and independent binomial trials. It uses only the interval rule: it does not model a task with broken tests being left out of M3, and it judges each case separately, whereas real forms are made at release scope and apply to every case alike. The assumption that defects only inflate scores does not hold for tests that reject a correct fix, which lower resolve rates. Tasks from one repository may be correlated, and agent runs vary between attempts; the Wilson interval does not account for either. The replay changes records in designed ways; their effect follows from the rules, so only the behavior on collected records is informative. Eligibility checks may remove tasks with unstable tests and so change the sampled workload. The exploit attempts are one attempt per task with at most three patches, so a failure to find a defect does not show that none exists. ABC's decision equals the resolve rate alone, apart from its caveats, because ABC has no rule to accept or reject a result, and the model's answers to the ABC items are not verified by a second person. Four limits concern the people and models. The author is the only signer of the forms, so no second person checks a signature. Claude Sonnet 5.5, which writes the rating forms and answers the ABC items, is the model the MESA paper used as its reviewer and is of the same family as Claude Opus 5.5, which built the engine, so agreement between them is not independent. The two rating runs share the model, the prompt, and the inputs; while run A worked, a concurrent rating attempt shared its scratch folder, and run A saw a few lines of that attempt's script, so run B was run alone afterwards. The ABC run also read the MESA rating prompt's rules, though not the rubric or any rating form. Claude Sonnet 5.5 also made the exploit attempts and writes the forms that cite them. Two runs of the same model agree with each other more than two people would.

\textbf{Construct validity.} Decisions depend on the seven selected criteria and on the policy; another policy or another choice of criteria would give other decisions, and each criterion added makes the strict policy harder to meet. A passed check within stated detection limits does not prove the absence of defects. Under the strict policy, abstention may be certain once the forms are signed, because a form that records no performed check for B12, B14a, or B14b cannot lead to a decision. The requirement is one acceptance criterion with illustrative values; resolved means that the maintainers' tests pass, which is narrower than a fix the maintainers would merge, and issues without tests are not covered. It does not address maintainability, security, or review effort. A difference between source and target results cannot be split into training exposure and task difficulty, because the source tasks are older and screened and the target tasks are newer and not screened; the interval comparison cannot separate the two. B14b is judged for each coding agent from public material, and the held out target window rests on the providers' stated cutoffs. Invalid sessions are detected from logs, so the count by kind is a lower bound.

\textbf{External validity.} One repository, one reviewing model, two coding agents, and 20 source and 20 target tasks per agent limit generalization; two coding agents give two cases for each question, not a rate; an interval around 20 tasks is still wide, and small samples leave most decisions inconclusive. Training data exposure (B14a, B14b) can rarely be verified for closed models from public material, which may lead to uniform abstention for both coding agents.

\textbf{Roles of the models and the author.} Claude Code with Claude Opus 5.5 implemented the records, engine, replay definitions, session controls, harness and session tooling, run scripts, templates, and tests; ran the eligibility checks, the target retrieval, the checks on the published environments, and the comparison with ABA's published findings; ran the sandbox checks and the rehearsal, chose the Antigravity settings, its hook, and the shared Python, and updated the Codex tool and locked both tools against updates at the author's request; drafted the requirement and the transfer justifications; launches the pilot sessions at the author's request; and revised this manuscript. Claude Sonnet 5.5, in separate runs, makes the B07 exploit attempts and the B28a grading check, writes the rating forms twice (runs A and B), and answers the ABC items. AI models in coding tools prepared the simulation scripts, and critiques generated by a model (GPT-6.1 Sol through the Codex command line tool) were used as editing input. The author decided the design, the coding agents, the target window, and the rules; reads and signs the rating forms, the ABC record, the requirement, and the transfer justifications; and authorizes each phase. The model activities and the passing tests are not independent validation of the levels, of the method, or of the manuscript.

\section{Conclusion}
\label{sec:conclusion}
This paper applies the requirements traceability rules of MESA to a requirement on a coding agent, with versioned records and a decision engine that keeps input integrity, evidence status, and decision apart and explains each decision. The implementation makes explicit what a MESA level must record to count for a requirement: a level, a check with its outcome and kind, a scope that covers the coding agent and tasks, an evidence location, and detection limits. Calculations in Section~\ref{sec:completed} show that small samples leave most decisions inconclusive, and that any benefit of rating forms in simulation depends on the assumed accuracy. The prepared SymPy evaluation, with 20 source and 20 target tasks per coding agent and at most 90 sessions, will show how a review with MESA or ABC changes the decision beyond an automated audit, and whether a decision from the reviewed benchmark agrees with a direct test of the requirement on new bug reports. Until it is conducted, the method's effect on requirement decisions remains untested.

\section*{Data Availability}
The replication package holds the full development record: the decision engine and record schema, the protocol records, the exploit attempts, the rating forms of both runs and their agreement, the ABC record, the study manifest, the SymPy task rows, the target pool search and target rows, the eligibility results with harness log hashes (the harness logs and the raw session and sandbox logs, which hold personal data, are kept outside version control; their hashes and parsed summaries are included), the checks run by the study, result templates, readiness report, the simulation scripts and outputs, and the test suites. An anonymized copy of the replication package is provided for review.
% TODO(author): add the anonymized link.

\bibliographystyle{splncs04}
\bibliography{references}

\end{document}
